\documentclass[10pt,a4paper]{amsart}
\usepackage[margin=19mm]{geometry}
\usepackage{amsmath,amssymb,mathtools}
\usepackage[scr=rsfs,cal=euler]{mathalfa}
\usepackage{xfrac}
\usepackage[unicode,hidelinks,bookmarks=false]{hyperref}
\newcommand{\ie}{i.e.\ }
\newcommand{\cf}{cf.\ }
\newcommand{\resp}{respectively\ }
\newcommand{\Subsection}{\subsection}
\providecommand{\nobreakdash}{\nobreak\hbox{-}}
\setlength{\emergencystretch}{2em}
\multlinegap0pt
\usepackage{amssymb}
\usepackage{mathtools}
\usepackage{float}
\usepackage[shortlabels]{enumitem}
\usepackage{tikz}
\newtheorem{lemma}{Lemma}
\newcommand{\Mapo}[2]{\underline{\mathrm{Map}}^{0}\left (#1,#2\right )}
\newcommand{\Pic}{\underline{\mathrm{Pic}}}
\newcommand{\Spec}{\mathrm{S}\mathrm{pec}\,}
\newcommand{\bA}{\mathbb{A}}
\newcommand{\bG}{\mathbb{G}}
\newcommand{\bP}{\mathbb{P}}
\newcommand{\cO}{\mathcal{O}}
\title{Equivariant Elliptic Cohomology and Mapping Stacks}
\author{Nicol\`o Sibilla, Paolo Tomasini}
\date{Source: arXiv:2303.10146v3 (CC BY 4.0)}
\begin{document}
\maketitle

\noindent\emph{Source and licence.} Equivariant Elliptic Cohomology and Mapping Stacks. Version arXiv:2303.10146v3 (CC BY 4.0), distributed by arXiv under the Creative Commons Attribution 4.0 International licence, http://creativecommons.org/licenses/by/4.0/, as recorded in the arXiv licence metadata for this exact version and shown on the abstract page https://arxiv.org/abs/2303.10146v3. Full licence notice: \texttt{licenses/arxiv-2303.10146-CC-BY-4.0.txt}.\par\smallskip

\noindent\textit{Editorial scope.} Counted unit: the article's lemma lemma:Pncodescent (source0.tex lines 1596--1598). The article's complete original proof of it is delivered with it (source0.tex lines 1599--1652, one \texttt{proof} environment of 54 lines). No mathematics is written, repaired or re-derived: the statement and the proof are the article's own lines, in the article's own notation, and no retained line is substituted or re-flowed.  Retained uncounted as the source-local context the proof needs: lines 1548--1551; lines 1587--1590.  Nothing was omitted from any retained span, so the package carries no omission record.  No retained line contains a citation, so the wrapper carries no bibliography and no bibliography entry.  No retained line contains a figure environment, an image-inclusion command or drawing code, so no image is delivered and the package declares no asset dependency.  Editorial formatting only, in addition: an AMS \texttt{amsart} preamble that restates the article's own declarations and macros used by the retained lines, namely the environments \texttt{lemma} on separate counters, together with the standard packages those lines need.  The retained environments print in this document as Lemma 1 in order of appearance, whereas the article prints the same items with the numbering of its own full document.  The complete native source of the article as distributed in the arXiv e-print for this exact version is bundled unchanged as \texttt{sources/138824/source0.tex}.
	\begin{lemma}
    \label{lemma:Ancodescent}
The stack of quasi-constant maps $\Mapo{E}{[\mathbb{A}^N/T]}$ satisfies $T$-equivariant Zariski codescent on $[\mathbb{A}^N/T]$. 
\end{lemma}

\begin{equation}
\label{standardform}
\lambda  \cdot [z_0, z_1 \ldots z_N] =    [z_0, \prod_{i=1}^{n}\lambda_i^{w^1_i}z_1,\dots,\prod_{i=1}^{n}\lambda_i^{w^N_i}z_N] \in \bP^N  
\end{equation} 

\begin{lemma}\label{lemma:Pncodescent}
The stack of quasi-constant maps $\Mapo{E}{[\mathbb{P}^N/T]}$ satisfies $T$-equivariant Zariski codescent on $[\mathbb{P}^N/T]$.
\end{lemma}

\begin{proof}
The proof strategy is the same as for Lemma \ref{lemma:Ancodescent}. Namely, consider a map $$f:E_K \to [\bP^N/T]$$  classified by a point of 
$\Mapo{E}{[\bP^N/T]}$. 
We need to show that $f$ factors through the image of a $T$-orbit in 
$[\mathbb{P}^N/T]$. Now we make the following observations:
\begin{enumerate}
\item all line bundles on $\mathbb{P}^N$ admit a $T$-equivariant structure;
\item there exists a $r > 0$ such that there exists a non-vanishing $T$-equivariant section $\sigma$ of $\cO_{\mathbb{P}^N}(r)$;
\item we can further assume that, in point $(2)$, $r=1$ on condition of replacing, if needed, $\mathbb{P}^N$ with a larger projective space $\mathbb{P}^M$ equipped with a $T$-action and such that there is a $T$-equivariant Veronese embedding 
$$
\mathbb{P}^N  \to \mathbb{P}^M
$$
\end{enumerate}
Although these are all standard facts, let us sketch a proof. We start with $(1)$. We identify $\mathbb{P}^N$ with the projectivization $\mathbb{P}(V)$ of the vector space $V$ which is  dual to $H^0(\bP^N,\cO_{\bP^N}(1))$. Note that we can lift the $T$-action on $\bP^N$ to a linear action on $V$. This turns $V$ into a $T$-representation. Taking the dual representation gives a  $T$-action on $H^0(\bP^N,\cO_{\bP^N}(1))$.  This induces a $T$-equivariant structure on the line bundle $\cO_{\bP^N}(1)$. Taking tensor powers and duals generates $T$-equivariant structures on all the bundles $\cO_{\bP^N}(m)$, for all integers $m$.

Let us consider $(2)$ next. The existence of such a section $\sigma$ for some $r$ is equivalent to fact that the GIT quotient $\bP^N//T$ is non-empty, i.e. it is equivalent to the existence of a semistable point of $\bP^N$ with respect to the $T$-action. A semistable point in $\bP^N$ with respect to a linear action is by definition a semistable point of the lift of the action to the vector space $V \cong k^{N+1}$ such that 
$\mathbb{P}^N \cong \mathbb{P}(V)$, i.e. a point in $V$ such that the closure of its orbit under the $T$-action does not contain $0$. The existence of a semistable point is clear after we write the action in standard form (see equation (\ref{standardform}) above). Indeed, the point $(1,0,0,\dots,0)$ in $k^{N+1}$ is fixed by the lift of the action, and hence it is in particular  semistable, as the closure of its orbit does not contain the point $0$. This implies that the point $[1:0:0:\dots:0]$ in $\bP^N$ is also semistable, hence the GIT quotient $\bP^N//T$ is nonempty. 

As for point $(3)$, i.e. the reduction to the case $r=1$, it is sufficient 
to linearize the action of $T$ on $\bP^N$ by choosing an equivariant structure on $\cO_{\bP^N}(r)$ which makes $\sigma$ into an equivariant section. As familiar from GIT theory, the choice of linearization yields a $T$-equivariant embedding
$$\bP^N \to \bP(H^0(\cO_{\bP^N}(r))^\vee) \cong \bP^M$$ 
In particular, we obtain an isomorphism of $T$-modules
\begin{equation*}
    H^0(\bP^M,\cO_{\bP^M}(1)) \cong H^0(\bP^N,\cO_{\bP^N}(r))
\end{equation*}
which restricts to  the spaces of equivariant sections
\begin{equation*}
    H^0(\bP^M,\cO_{\bP^M}(1))^{T} \cong H^0(\bP^N,\cO_{\bP^N}(r))^{T}
\end{equation*}
hence a bijective correspondence between the $T$-equivariant sections of the two bundles. Note that $f$ factors through the image of a $T$-orbit if and only the composite map
$$
E_K \to [\bP^N/T] \to [\bP^M/T]
$$ 
factors through the image of a $T$-orbit. Thus, from the perspective of the argument we are carrying out, we can harmlessly replace $\bP^N$ with $\bP^M$.  We do this implicitly in the sequel, and in particular assume that $r=1$ and $\sigma$ is linear, but we will not rename either $\bP^N$ or $\sigma$. 



Consider the map $\phi$  classifying the line bundle $\cO_{[\bP^N/T]}(r)$, $\, 
\phi: [\bP^N/T] \to [\Spec k/\bG_m].$   
%Note that composition yields a map
%\begin{equation}
%\label{sec0}
%\Mapo{E}{[\bP^N/T]} \to  \Mapo{E}{[\ast/\bG_m]} \cong \Pic^0(E)
%\end{equation}
If the map $f$ is classified by a geometric point  of $\Mapo{E}{[\bP^N/T]}$, the pullback line bundle $f^*(\cO_{[\bP^N/T]}(r))$ is classified by the composition 
$$
\Spec(K) \stackrel{f}\to \Mapo{E}{[\bP^N/T]} \stackrel{\phi}\to \Mapo{E}{[\Spec k/\bG_m]} \cong \Pic^0(E)
$$
In particular, $f^\ast(\cO_{[\bP^N/T]}(r))$ is a degree-zero line bundle on $E$. 

Let $H \subset \bP^N$ be the zero locus of $\sigma$, and consider its complement $U$. As $\sigma$ is linear $H$ is a hyperplane and $U$ is a $T$-equivariant open subset of $\bP^N$ isomorphic to $\mathbb{A}^N$. We claim that the image of $f$ is entirely contained either in $H$ or in $U$. Indeed, assume to the contrary that the image of  $f$ intersects both $H$ and $U$. The pullback section $f^\ast\sigma \in H^0(E, f^\ast\cO_{[\bP^N/T]}(r))$ is not constant, as the image of $f$ intersects both the zero-locus of $\sigma$ and its complement. However the line bundle $f^\ast\cO_{[\bP^N/T]}(r)$ must be of degree zero, and therefore its sections are necessarily constant. Thus as we claimed $f$ factors either through 
$$[\bA^N/T] \cong [U/T] \subset [\bP^N/T] \quad \text{ or through} \quad [H/T] \subset [\bP^N/T]$$ 
In the first case, we are done by Lemma \ref{lemma:Ancodescent}. In the second case, $f$ factors through the lower dimensional  space $[H/T] \cong [\bP^{N-1}/T]$, and we can conclude by induction on the dimension $N$: note indeed  that the base case of the statement $N=1$ is clear, as the complement of a $T$-invariant open subset $\bA^1 \cong U \subset \bP^1$ is a fixed point, which is a $T$-orbit. 
\end{proof}

\end{document}
